% ---------------------------------------------------------------------
\section{前向 Compton 振幅：解析性、交叉与色散关系（部分子靶）}\label{sec:dis-compton}

$W^{\mu\nu}$ 由对易子定义，只在物理区 $0<x\le1$ 有意义；OPE 却是关于编时乘积在欧氏动量下的陈述。连接二者的是前向虚 Compton 振幅：它的吸收部分是 $W^{\mu\nu}$，它的 Taylor 系数在非物理点 $\omega=0$ 处计算。本节对部分子靶建立这一联系；强子靶的对应陈述需要另外的工具，放在\cref{sec:dis-np}。

\begin{definition}[前向 Compton 振幅]\label{def:dis-compton}
\begin{equation}\label{eq:dis-T}
T^{\mu\nu}(P,q)=\ii\int\dd^4z\,\ee^{\ii q\cdot z}\langle P|T\,j^\mu(z)j^\nu(0)|P\rangle
=-\tilde g^{\mu\nu}\,T_1(\omega,Q^2)+\frac{\tilde P^\mu\tilde P^\nu}{P\cdot q}\,T_2(\omega,Q^2),\qquad \tilde g^{\mu\nu}=g^{\mu\nu}-\frac{q^\mu q^\nu}{q^2},
\end{equation}
$\tilde P$ 与\cref{eq:dis-W}相同；$\omega=2P\cdot q/Q^2$。在 $\omega=0$ 附近，$T_2$ 的 Taylor 展开记作 $T_2=\sum_{n}t_n\,\omega^{n-1}$（$t_n$ 称为 $T_2$ 的 Taylor 系数），$T_1=T_1(0)+\sum_n\tau_n\,\omega^n$。% VF-CONV-14
对强子靶，编时乘积在 $z=0$ 处的协变定义只确定到接触项（局域多项式），这一点由\cref{asm:intro-wightman}(W4) 处理；对部分子靶，$T^{\mu\nu}$ 就是截腿的前向 Feynman 图之和。
\uses{def:dis-kinematics}
\end{definition}

领头阶的两个图（$s$ 道与 $u$ 道“手提包”图）见\cref{fig:forward-compton}。

%%INSERT fig:forward-compton

\begin{proposition}[解析性与交叉（部分子靶） \tagR]\label{prop:dis-compton-analytic}
（$R_{\mathrm{PT}}$：在重整化微扰论中对红外正规化的部分子态严格。）设靶是夸克或胶子，其红外发散由外腿的类空离壳性 $p^2=-P^2<0$ 正规化（内线质量任意，$\ge0$）。$T^{\mu\nu}$ 为截腿的前向 Green 函数对外腿自旋（极化）的平均，$W^{\mu\nu}$ 为相应割图之和；$T_1$、$T_2$（以及 $F_1$、$F_2$）定义为 $T^{\mu\nu}$（以及 $W^{\mu\nu}$）的张量分解中 $-g^{\mu\nu}$ 的系数与 $p^\mu p^\nu/\nu$ 的系数（$T^{\mu\nu}$ 横向时即\cref{eq:dis-T}中的不变振幅）。在微扰论的每一阶、固定 $Q^2>0$，记 $\omega_{\mathrm{th}}=1+P^2/Q^2$：
\begin{enumerate}[label=(\alph*),leftmargin=2.2em,itemsep=0pt,topsep=2pt]
\item $T_1$、$T_2$ 在 $\mathcal C_{\mathrm{th}}=\mathbb C\setminus\big((-\infty,-\omega_{\mathrm{th}}]\cup[\omega_{\mathrm{th}},\infty)\big)$ 中解析（因而在沿实轴 $|\omega|\ge1$ 割开的平面中解析），并且是实解析的：$T_i(\omega^*)=T_i(\omega)^*$；
\item $T_1(-\omega)=T_1(\omega)$，$T_2(-\omega)=-T_2(\omega)$；
\item 对实的 $\omega\ge1$，作为分布 $\Imag T_i(\omega+\ii0)=2\pi F_i(1/\omega)$（$1\le\omega<\omega_{\mathrm{th}}$ 时两边都为零）。
\end{enumerate}
（本命题只涉及部分子态。）
\uses{def:dis-compton,def:ir-landau-setup,prop:prelim-optical}
\end{proposition}

\begin{proof}
(b) 前向图的集合在交换两个流顶点 $(\mu,q)\leftrightarrow(\nu,-q)$ 下不变，故 $T^{\mu\nu}(p,q)=T^{\nu\mu}(p,-q)$。在 $q\to-q$ 下 $\nu=p\cdot q\to-\nu$，$\omega\to-\omega$。自旋（极化）平均且不含 $\gamma_5$ 的前向振幅关于 $\mu\nu$ 对称：反对称结构只可能是 $p^\mu q^\nu-q^\mu p^\nu$ 或 $\epsilon^{\mu\nu\rho\sigma}p_\rho q_\sigma$，后者需要 $\gamma_5$ 或自旋，前者与对称的 $g^{\mu\nu}$、$p^\mu p^\nu$ 的系数无关。于是 $g^{\mu\nu}$ 的系数是 $\omega$ 的偶函数，$p^\mu p^\nu$ 的系数也是偶函数，乘以 $\nu$ 后 $T_2$ 为奇函数。

(c) 由割图规则，每一阶 $2\Imag T^{\mu\nu}$ 等于前向图在 $s$ 道（$P\cdot q>0$）的全部割之和，即 $4\pi W^{\mu\nu}$（\cref{prop:prelim-optical}(b)；对离壳外腿这是关于 Feynman 积分的恒等式，$W^{\mu\nu}$ 即按此定义）。取张量系数即得 (c)；$u$ 道的割要求 $(p-q)^2>0$，在 $\omega>0$ 时不出现。

(a) 张量积分在 Feynman 参数（$\alpha$）表示中直接给出各张量结构 $g^{\mu\nu}$、$p^\mu p^\nu$、$p^\mu q^\nu$、$q^\mu p^\nu$、$q^\mu q^\nu$ 的系数：圈动量平移到 $k_c(\alpha)$（$p$、$q$ 的线性组合，系数是 $\alpha$ 的有理函数）后作 Gauss 积分，得到形如
\begin{equation}\label{eq:dis-alpha-rep}
\int_{\Delta}\dd\mu(\alpha)\;h(\alpha,\omega)\,\mathcal U(\alpha)^{a}\,\mathcal F(\alpha,\omega)^{-c}
\end{equation}
的参数积分（\cref{def:ir-landau-setup}），其中 $h$ 是 $\alpha$ 与 $\omega$ 的多项式，$a,c$ 依赖于 $D$；不出现 Gram 行列式。第二 Symanzik 多项式是对图的生成 2-森林 $T$ 的求和，$\mathcal F=\sum_T\alpha^{E\setminus T}(-s_T)+\mathcal U\sum_j\alpha_jm_j^2$，$s_T$ 是流入 $T$ 的一个连通分支的外动量之和的平方。前向运动学中流入一个分支的外动量为 $c_pp+c_qq$，$c_p,c_q\in\{0,\pm1\}$，故
\[
-s_T=c_p^2P^2+c_q^2Q^2-c_pc_q\,\omega Q^2,\qquad \mathcal F=A(\alpha)+\omega B(\alpha),
\]
$A=\sum_T\alpha^{E\setminus T}(c_p^2P^2+c_q^2Q^2)+\mathcal U\sum_j\alpha_jm_j^2\ge0$，$B=-Q^2\sum_Tc_pc_q\,\alpha^{E\setminus T}$ 为实数。因为两个流顶点不同，存在分开它们的 2-森林（$c_q\neq0$），所以在单纯形 $\Delta$ 的内部 $A>0$。对 $c_pc_q\ne0$ 的单项式，$A$ 中的系数是 $P^2+Q^2=Q^2\omega_{\mathrm{th}}$，而 $|B|$ 中的系数是 $Q^2$。由此对 $\omega=\omega_1+\ii\omega_2$：若 $|\omega_1|<\omega_{\mathrm{th}}$，则 $\Real\mathcal F=A+\omega_1B\ge(1-|\omega_1|/\omega_{\mathrm{th}})A>0$；若 $|\omega_1|\ge\omega_{\mathrm{th}}$ 而 $\omega_2\ne0$，则 $|\mathcal F|^2=(A+\omega_1B)^2+\omega_2^2B^2\ge(1-|\cos\theta|)(A^2+|\omega|^2B^2)$（$\theta=\arg\omega$），故 $|\mathcal F|\ge\sqrt{1-|\cos\theta|}\,A$。总之在 $\mathcal C_{\mathrm{th}}$ 的每个紧子集 $K$ 上 $|\mathcal F|\ge c_KA$，$c_K>0$。于是\cref{eq:dis-alpha-rep}的被积函数在 $K$ 上被 $\omega=0$ 处的被积函数乘以常数所控制；在参数积分绝对收敛的 $D$ 区域内，由控制收敛与 Morera 定理，积分在 $\mathcal C_{\mathrm{th}}$ 中解析。到 $D=4-2\varepsilon$ 附近的亚纯延拓也保持这一解析性：$B$ 的每个单项式都出现在 $A$ 中，故对 $\omega\in\mathcal C_{\mathrm{th}}$，$\mathcal F$ 与 $A$ 的 Newton 多面体相同，扇区分解在每个扇区中把二者写成同一个单项式乘以满足同样下界的非零因子；$\overline{\rm MS}$ 抵消项是外动量的多项式，因而是 $\omega$ 的多项式。$A$、$B$ 为实数给出实解析性。$T_1=-(g^{\mu\nu}$ 的系数$)$，$T_2=\nu\times(p^\mu p^\nu$ 的系数$)$ 都是这样的函数的有限和。
\end{proof}

对外腿在壳的部分子（例如在纯维数正规化中的无质量夸克与胶子），流守恒使 $T^{\mu\nu}$ 横向，上面的系数就是\cref{eq:dis-T}中的不变振幅；那里相应的结论是逐个 Laurent 系数地证明的，并且胶子靶在 $\omega=0$ 的正则性只有全阶论证（附录~M 的\cref{cor:appM-cut-plane}）。强子靶的解析性需要 Wightman 公理，见\cref{prop:dis-dispersion-np}。

\cref{fig:omega-plane}画出了 $\omega$ 平面的解析结构与下面证明所用的围道。

%%INSERT fig:omega-plane

\begin{proposition}[色散关系：Taylor 系数 = 矩（部分子靶） \tagR]\label{prop:dis-dispersion}
（$R_{\mathrm{PT}}$：在重整化微扰论中对红外正规化的部分子态严格。）在\cref{prop:dis-compton-analytic}的设定下固定微扰论的阶。设在这一阶满足\emph{部分子增长条件}：当 $|\omega|\to\infty$ 时，在闭割平面上（含割线上的边界值）一致地有
\begin{equation}\label{eq:dis-pt-growth}
T_2(\omega)=o(|\omega|),\qquad T_1(\omega)=o(|\omega|^2).
\end{equation}
则对 $|\omega|<1$，
\begin{equation}\label{eq:dis-dispersion}
T_2(\omega)=\sum_{n=2,4,\dots}4\,\omega^{n-1}M^{(2)}_n,\qquad T_1(\omega)-T_1(0)=\sum_{n=2,4,\dots}4\,\omega^nM^{(1)}_n,
\end{equation}
\begin{equation}\label{eq:dis-moments-def}
M^{(2)}_n=\int_0^1\dd x\,x^{n-2}F_2(x),\qquad M^{(1)}_n=\int_0^1\dd x\,x^{n-1}F_1(x)
\end{equation}
（含弹性峰；积分收敛）。树图阶的增长条件成立。% VF-CONV-14
（本命题只涉及部分子态。）
\uses{prop:dis-compton-analytic,def:dis-compton}
\end{proposition}

\begin{proof}
设 $n$ 为偶数。由 Cauchy 公式，$t_n=\frac1{2\pi\ii}\oint_{C_0}T_2(\omega)\,\omega^{-n}\,\dd\omega$，$C_0$ 是原点附近的小圆（\cref{fig:omega-plane}(a)）。由\cref{prop:dis-compton-analytic}(a)，把 $C_0$ 变形为半径 $R$ 的大圆加上紧贴两条割线 $\pm[1,R]$ 的发夹形路径（\cref{fig:omega-plane}(b)；$R$ 取在边界值连续的点）。沿右割线的贡献为
\[
\frac1{2\pi\ii}\int_1^R\Big[T_2(\omega+\ii0)-T_2(\omega-\ii0)\Big]\omega^{-n}\dd\omega
=\frac1{2\pi\ii}\int_1^R2\ii\cdot2\pi F_2(1/\omega)\,\omega^{-n}\dd\omega=2\int_{1/R}^1\dd x\,x^{n-2}F_2(x),
\]
这里用了实解析性 $T_2(\omega-\ii0)=T_2(\omega+\ii0)^*$、\cref{prop:dis-compton-analytic}(c) 与 $\omega=1/x$（$\omega^{-n}\dd\omega=-x^{n-2}\dd x$）。左割线上令 $\omega=-\omega'$：由奇性，$T_2(-\omega'\pm\ii0)=-T_2(\omega'\mp\ii0)$，故左割线上的跳跃等于右割线在 $\omega'$ 处的跳跃，而 $\omega^{-n}=\omega'^{-n}$（$n$ 偶），定向的反转与 $\dd\omega=-\dd\omega'$ 抵消，左割线给出同样的 $2\int_{1/R}^1x^{n-2}F_2$。大圆的贡献不超过 $R^{1-n}\sup_{|\omega|=R}|T_2|$，由\cref{eq:dis-pt-growth}当 $R\to\infty$ 时趋于零（$n\ge2$）。于是 $t_n=4M^{(2)}_n$，且 $\int_{1/R}^1x^{n-2}F_2$ 的极限存在。对奇数 $n$，$T_2$ 的奇性使 $\omega^{n-1}$ 的系数为零。对 $T_1$ 与偶数 $n\ge2$ 作同样计算：$\tau_n=\frac1{2\pi\ii}\oint T_1\omega^{-n-1}\dd\omega$，右割线给出 $2\int x^{n-1}F_1$；$T_1$ 是偶函数，左割线上的跳跃变号，而 $\omega^{-n-1}$ 也变号，两者抵消，再得 $2\int x^{n-1}F_1$；大圆贡献不超过 $R^{-n}\sup|T_1|\to0$。$n=0$ 时大圆一般不趋于零，$T_1(0)$ 因而是一个减除常数。

树图阶：\cref{ex:dis-free-quark}给出 $T_2=4e_q^2\omega/(1-\omega^2)=\order{1/\omega}$，$T_1=2e_q^2\omega^2/(1-\omega^2)=\order{1}$；外腿离壳时两条传播子分母为 $(p+q)^2=Q^2(\omega-\omega_{\mathrm{th}})$、$(p-q)^2=-Q^2(\omega+\omega_{\mathrm{th}})$，同样的计算给出 $T_2=4e_q^2\omega\,\omega_{\mathrm{th}}/(\omega_{\mathrm{th}}^2-\omega^2)$ 与有界的 $T_1$：极点从 $\pm1$ 移到 $\pm\omega_{\mathrm{th}}$，$|\omega|\to\infty$ 的行为不变。
\end{proof}

\begin{remark}[部分子增长条件的地位]\label{rem:dis-pt-growth}
(i) 固定阶的振幅在割平面上多项式有界：在参数积分收敛的区域中，上面证明里的 $|\mathcal F|\ge\sqrt{1-|\cos\theta|}\,A$ 直接给出 $|T_2(\omega)|\le C(1+|\omega|)^{A}|\Imag\omega|^{-B}$，$A,B$ 依赖于阶。这样的界只容许带有限个减除的色散关系（对强子靶的证明见\cref{prop:dis-dispersion-np}，那里还用到 $F_2\ge0$），它本身不能排除最低几个矩上的“固定极点”，即 $T_2$ 中不被割线跳跃决定的 $\omega$ 的奇多项式项；\cref{prop:dis-low-moments}对强子作了同样的分析。(ii) \cref{eq:dis-pt-growth}在树图阶成立。在单圈阶，对纯维数正规化中在壳的无质量部分子，前向振幅硬区展开的 Taylor 系数与割图结果的 Mellin 矩对 $N\le8$ 作为 $\varepsilon$ 的有理函数精确相等（\cref{ex:dis-oneloop-ope}），（硬区展开的系数就是在壳振幅的 Taylor 系数，\cref{thm:appM-forward-expansion}），即这些 $N$ 上没有固定极点；单圈的部分子结构函数 $\hat F_2=z\,C_2(z)$ 在 $z\to0$ 时为 $z\ln(1/z)$ 型（单圈系数函数在 $z\to0$ 时只有对数增长，见第~10~章的\cref{prop:disall-nlo}），也与\cref{eq:dis-pt-growth}一致。(iii) 在一般的阶，按 Regge 极限的幂次计数（$t$ 道交换的最高自旋是胶子的 $1$），预期 $T_2\sim\omega^0\ln^k\omega$、$T_1\sim\omega\ln^k\omega$，正好满足\cref{eq:dis-pt-growth}；但对每一阶、每个图（以及规范不变的图之和）的证明并不存在。附录~M 对在壳无质量部分子作了同样的判断：矩恒等式只有全阶论证（\cref{cor:appM-cut-plane}，\tagA）。强子靶的对应条件是\cref{asm:intro-growth}。
\end{remark}

\begin{example}[自由夸克的前向振幅]\label{ex:dis-free-quark}
取在壳无质量夸克靶（电荷 $e_q$），树图阶。按\cref{def:dis-compton}，$\psi(z)$ 与 $\bar\psi(0)$ 缩并的项给出 $s$ 道图，$\psi(0)$ 与 $\bar\psi(z)$ 缩并的项给出 $u$ 道图（\cref{fig:forward-compton}）。用 $\ii\int\dd^4z\,\ee^{\ii(q+p)\cdot z}S_F(z)=\ii\cdot\frac{\ii(\sla p+\sla q)}{(p+q)^2+\ii0}$ 并对自旋与颜色平均（$\frac12\sum_s\bar u\Gamma u=\frac12\TrD[\sla p\Gamma]$），
\[
T^{\mu\nu}=-e_q^2\,\frac{\frac12\TrD[\sla p\gamma^\mu(\sla p+\sla q)\gamma^\nu]}{(p+q)^2+\ii0}-e_q^2\,\frac{\frac12\TrD[\sla p\gamma^\nu(\sla p-\sla q)\gamma^\mu]}{(p-q)^2+\ii0},
\qquad \tfrac12\TrD[\sla p\gamma^\mu\sla k\gamma^\nu]=2\big(p^\mu k^\nu+p^\nu k^\mu-g^{\mu\nu}p\cdot k\big),
\]
在任意 $D$ 中成立。由 $p^2=0$ 与 $(p\pm q)^2=\pm2p\cdot q-Q^2=-Q^2(1\mp\omega)$，把张量写成 $g$、$pp$、$pq$、$qp$、$qq$ 的组合并与\cref{eq:dis-T}比较（$q_\mu T^{\mu\nu}=0$ 可直接验证），得
\begin{equation}\label{eq:dis-free-T}
T_2=\frac{4e_q^2\,\omega}{1-\omega^2}=2e_q^2\Big[\frac1{1-\omega}-\frac1{1+\omega}\Big],\qquad T_1=\frac{2e_q^2\,\omega^2}{1-\omega^2}
\end{equation}
（\texttt{s2\_compton.py} 精确到 $D$ 核对）。它们满足 $T_1$ 偶、$T_2$ 奇，$T_1(0)=0$，Taylor 展开 $T_2=\sum_{n=2,4,\dots}4e_q^2\omega^{n-1}$，$T_1=\sum_{n=2,4,\dots}2e_q^2\omega^n$。

物理区的 $\ii0$ 来自 $s$ 道传播子：$\frac{-1}{(p+q)^2+\ii0}=\frac{-1}{Q^2(\omega-1)+\ii0}$，故 $\Imag\frac{-1}{Q^2(\omega-1)+\ii0}=\pi\delta\big(Q^2(\omega-1)\big)>0$，这对应 $\omega\to\omega+\ii0$。于是由 $\frac1{1-\omega-\ii0}=\mathrm P\frac1{1-\omega}+\ii\pi\delta(1-\omega)$，
\[
\Imag T_2(\omega+\ii0)=2\pi e_q^2\,\delta(\omega-1)=2\pi F_2(1/\omega),\qquad F_2(x)=e_q^2\,\delta(1-x)
\]
（$\delta(1-1/\omega)=\omega^2\delta(\omega-1)=\delta(\omega-1)$）；同样 $F_1=\frac12e_q^2\delta(1-x)$，满足 $F_2=2xF_1$（树图 $F_L=0$）。矩为 $M^{(2)}_n=e_q^2$，$M^{(1)}_n=\frac12e_q^2$，与 Taylor 系数 $t_n=4e_q^2$、$\tau_n=2e_q^2$ 恰好满足\cref{eq:dis-dispersion}；而 $|\omega|\to\infty$ 时 $T_2\to0$、$T_1\to-2e_q^2$：$T_1$ 的减除常数 $T_1(0)$ 是必要的。$1/(4\pi)$ 的归一化正好给出部分子模型的 $F_2=e_q^2x\,\delta(1-x)$。
\uses{prop:dis-dispersion}
\end{example}
